\anglais

\cleardoublepage
\phantomsection
\addcontentsline{toc}{chapter}{General Conclusion}
\label{chap:general_conclusion}
\markboth{General Conclusion}{General Conclusion}

{\LARGE\bfseries General Conclusion\par}
\vspace{2\baselineskip}


This thesis studied solar flare prediction from active-region magnetograms through the lens of input design. The central question was not whether a new neural-network architecture could be proposed, but rather how much of the forecasting performance of modern deep-learning models depends on the way the solar data are represented before they reach the model. This distinction is important. In many computer-vision problems, the input format is treated as a practical preprocessing detail. In solar flare forecasting, however, the input representation determines which physical information is available: the spatial scale of the magnetic field, the position of the active region on the solar disk, the duration of the observed history, and the density of temporal sampling.

The work presented in this thesis therefore approached flare prediction as both a machine-learning problem and a data-design problem. Using the reduced-resolution active-region magnetogram dataset of Boucheron et al.~\cite{ref_boucheron}, the experiments in Chapter~\ref{chap:article1} evaluated four practical choices under a controlled fine-tuning protocol: image resolution, solar-disk position, temporal cadence, and number of frames. The study focused on two prediction tasks, $\geq C$ and $\geq M$ flares in the following 24 hours. This separation was essential because the two tasks behave differently. The $\geq C$ task provides a broader measure of flare productivity, while the $\geq M$ task is closer to the rare-event regime that is more relevant for high-impact space-weather forecasting but also more sensitive to class imbalance, threshold selection, and seed variability.

The first main conclusion is that spatial detail matters, but the loss of predictive skill with lower resolution is not abrupt. Full $224\times224$ magnetograms gave the strongest average performance in the image-based study, particularly for the rarer $\geq M$ task. At the same time, downsampled inputs at $112\times112$ and $56\times56$ retained a substantial fraction of the performance. This suggests that the model does not rely only on extremely fine pixel-scale structures. Large-scale polarity organization, active-region size, and broad magnetic complexity already carry meaningful predictive information. The advantage of full resolution, especially for stronger flares, nevertheless indicates that finer magnetic gradients or more localized structures may still contribute to the classification of more severe events.

The second conclusion is that disk position should not be treated as a neutral variable. Active regions closer to disk center gave stronger performance than those at larger longitudes, and training on central regions improved performance on the central test subset in the experiments reported in Chapter~\ref{chap:article1}. This result is consistent with the physical and observational difficulty of line-of-sight magnetograms away from disk center: foreshortening and projection effects make the observed magnetic structure less directly comparable across positions. From a practical perspective, this means that a forecasting pipeline may benefit from explicitly accounting for disk position, either through specialized models, geometric corrections, vector-field information, or evaluation protocols that report performance by longitude range rather than only as a single global score.

The strongest empirical pattern in this thesis came from the temporal experiments. For fixed 16-frame sequences, increasing the cadence from 60 to 180 minutes increased the temporal span from 15 to 45 hours and led to the best mean performance in the augmented MViT sequence experiments. This indicates that, within the tested range, broader temporal context was more useful than dense short-term sampling. In other words, the model benefited more from observing active-region evolution over many hours than from seeing many closely spaced frames that may contain redundant information. This finding is consistent with the physical intuition that flare-relevant magnetic configurations develop over time and may not be fully captured by a single magnetogram.

The frame-count experiment refined this conclusion. When the temporal span was fixed at 45 hours, shorter sequences with 4 or 8 frames remained competitive with the full 16-frame sequence, although the best choice depended on the task and metric. This is an important practical result. It suggests that temporal coverage and temporal density should be treated as separate design variables. A forecasting system may be able to reduce memory use, data loading cost, and inference time by sampling a long temporal window sparsely, rather than processing every intermediate observation. For large-scale operational pipelines, this kind of input efficiency is not only computationally convenient; it can determine whether repeated retraining, uncertainty estimation, or multi-model evaluation is feasible.

Together, these results support the main claim of the thesis: input design is a central part of deep-learning-based solar flare prediction. Architecture matters, but the architecture can only learn from the physical and temporal cues that the input representation preserves. Resolution controls the spatial scale of the magnetic structures available to the model. Disk position controls the amount of observational distortion mixed into the data. Cadence and frame count control whether the model observes short-term redundancy or longer-term magnetic evolution. These choices shape the forecasting problem before optimization begins.

The thesis also emphasizes the importance of evaluation under class imbalance. Accuracy is not sufficient for this problem, especially for $\geq M$ prediction, because a high value can be obtained by overpredicting the majority non-flaring class. For this reason, the experiments relied on skill scores and ranking metrics such as TSS, ROC-AUC, PR-AUC, and F1 at a validation-selected threshold. The appendices further expand this view using additional operating-point metrics such as HSS and MCC, especially in Appendix~\ref{app:extended_image_operating_metrics} and Appendix~\ref{app:extended_temporal_metrics}. The broader lesson is that rare-event forecasting cannot be summarized by one number. A useful model must be assessed both as a ranking system and as a thresholded forecasting system, because good ranking performance does not automatically imply a stable operational decision threshold.

Several limitations should be kept in mind. First, the experiments were conducted on one benchmark dataset derived from Solar Cycle 24 and from one preprocessing pipeline. The conclusions are therefore controlled evidence about this dataset, not a complete operational validation across solar cycles, instruments, and real-time conditions. Second, the work used line-of-sight magnetograms represented as reduced-resolution 8-bit images. These inputs are convenient for computer-vision models and efficient for experimentation, but they do not contain the full vector magnetic-field information that may be important for flare physics. Third, the experiments were designed to isolate input factors rather than to exhaustively tune every model architecture, loss function, calibration method, or sampling strategy. Finally, the results are mostly quantitative without the model's reasoning. They show that some input representations contain useful predictive information, but they do not yet prove which physical structures the models use internally.

These limitations point directly to future work. A first direction is to decouple cadence and temporal span more systematically by evaluating a broader grid of temporal windows, frame counts, and sampling patterns. The results of this thesis suggest that long temporal context is valuable, but the optimal way to sample that context remains open. A second direction is to improve physical interpretability. Saliency maps, occlusion studies, attention analysis, and case-based studies could help determine whether the models focus on plausible flare-related structures, such as polarity inversion lines, compact high-gradient regions, or broader changes in magnetic complexity. A third direction is to extend the input modalities. Combining magnetogram images with scalar magnetic features, GOES X-ray histories, vector magnetograms, or other solar observations may provide complementary information that single-modality image models cannot capture.

Future work should also examine generalization more rigorously. Testing on later solar cycles, on data from different instruments, or under more operational near-real-time constraints would clarify whether the observed input-design trends remain stable outside the benchmark setting. Calibration and uncertainty estimation are also important for deployment: an operational forecast is not only a binary decision, but often a probability or risk estimate that must support human decision-making. Finally, dataset issues such as overlapping active regions, label ambiguity, and the dependence between nearby samples should be studied more carefully, since they can influence both training and evaluation.

In conclusion, this thesis shows that deep learning can extract meaningful predictive information from active-region magnetograms, but it also shows that the success of such models depends strongly on how the data are presented to them. The main contribution is therefore practical as well as methodological: lower spatial resolutions and sparse temporal sampling can remain competitive in some settings, disk-position effects should be explicitly considered, and longer temporal context appears to be one of the clearest drivers of sequence-model performance. These findings provide guidance for designing more efficient and better-controlled flare-prediction pipelines. More broadly, they support a view of solar flare forecasting in which model architecture, physical data representation, and evaluation protocol must be developed together rather than treated as separate components.
